\section*{如何使用本书}

这不是一本把公式按字典排列的手册。每章从一个量化研究中会真正遇到的问题出发，先说明为什么需要新的数学对象，再逐步建立定义、直觉和定理，随后用完整算例和反例检验结论的边界。读者应在每个主要推导处停下来，尝试解释每个假设的作用；如果某个假设在证明中没有出现，通常意味着命题可以变强，或者证明漏掉了关键一步。

上册的章节按照依赖关系排列，但不要求所有读者走同一条路线。希望尽快进入应用的读者，可以优先学习极限、线性代数、概率分布、统计推断、时间序列、优化、数值稳定性和 Monte Carlo；希望建立严格基础的读者，则从分析、测度和条件期望开始顺读。每章末尾的练习从概念解释逐渐过渡到证明、计算和研究判断，完整答案与常见误区收录在共享答案册中。

Notebook 展示数学对象如何一步步变成计算：先预测现象，再观察中间量与图形，最后修改参数并解释变化。正文中的小算例可用纸笔理解，复杂实验无需事先手算全部输出。解析结果帮助核对计算；改变数据后应重新推理，不必维持原来的小数。完整推导与练习反馈在共享答案册中。

本书只在这里解释一次阅读方式。项目维护者所需的文件布局、Notebook 命令、图形生成和出版流程另见仓库文档，不属于读者必须掌握的数学内容。
